\section{Feature Comparison}
\label{sec:feature-comparison}

Section~\ref{sec:taxonomy} introduced five classification axes and
applied them narratively to the QIEA-for-MOO methods identified so far.
This section applies the same axes systematically, one row per method,
across two tables: Table~\ref{tab:feature-comparison-taxonomy} recaps
representation, operator type, selection mechanism, and application
domain (the four axes for which the notes gathered while building
Section~\ref{sec:background}'s and Section~\ref{sec:taxonomy}'s
citations already give a confident answer), and
Table~\ref{tab:feature-comparison-performance} turns to complexity,
scalability, convergence behavior, benchmarks, and reported performance
-- the dimensions this survey's contributions
(Section~\ref{sec:introduction:contributions}) also promise to cover.

A methodological caveat applies to Table~\ref{tab:feature-comparison-performance}
specifically, and is itself a preview of
Section~\ref{sec:critical-analysis}'s reproducibility discussion: at
this drafting stage, several of its cells reflect only what each
method's title, abstract, and this survey's own literature-gathering
notes state, not an independent re-derivation or reproduction of the
original results. Cells that would require reading the full paper to
state with confidence are marked accordingly rather than filled with an
unverified number, following the same verify-before-citing discipline
already applied to every entry in this survey's bibliography. That so
many cells resist confident filling from title/abstract alone, for a
comparison table that is supposed to be the survey's most concrete
deliverable, is itself early evidence for a pattern
Section~\ref{sec:critical-analysis} treats as a cross-cutting limitation:
QIEA-MOO papers routinely under-report the complexity, scale, and
statistical rigor of their own evaluations relative to what a reader
would need to compare them fairly against one another.

\begin{table}[htbp]
  \centering
  \small
  \caption{Representation, operator, selection, and domain summary of
    the surveyed QIEA-MOO methods (recap of the taxonomy axes,
    Section~\ref{sec:taxonomy}).}
  \label{tab:feature-comparison-taxonomy}
  \begin{adjustwidth}{-2cm}{-2cm}
  \centering
  \begin{tabularx}{\linewidth}{@{}p{2.1cm}XXXX@{}}
    \toprule
    Method & Representation & Quantum-Inspired Operator(s) & Selection /
    Diversity & Application Domain \\
    \midrule
    QMEA~\citep{kim2006qmea} &
      Q-bit chromosome, observed to a binary decision vector &
      Rotation-gate update only &
      Direct Pareto-dominance ranking (no crowding distance) &
      Multi-objective 0/1 knapsack (combinatorial) \\
    QNSGA-II~\citep{guzel2022qnsga2} &
      Q-bit chromosome layered onto an NSGA-II individual &
      Rotation gate + superposition-based sampling &
      NSGA-II crowding distance, unmodified &
      General-purpose multi-objective benchmarks (not itemized in
      abstract-level notes) \\
    Continuous QIEA-MOO~\citep{olvera2023continuous} &
      Q-bit angle decoded to a real-valued decision variable &
      Rotation gate &
      Compared directly against NSGA-II as an external baseline &
      Continuous-space multi-objective benchmark problems \\
    EAH-QNSGA-II~\citep{kumar2026eahqnsga2} &
      Q-bit chromosome layered onto an NSGA-II individual &
      Adaptive rotation gate + entanglement-driven crossover +
      post-hoc local search &
      NSGA-II crowding distance, unmodified &
      Electric-vehicle charging-station siting (real-world
      infrastructure) \\
    QIEA-MOO for waste collection~\citep{gilea2026waste} &
      Not yet confirmed against full text (self-citation; venue/DOI
      itself unconfirmed, see the bibliography note) &
      Not yet confirmed against full text &
      Not yet confirmed against full text &
      Waste-collection vehicle routing (real-world logistics) \\
    \bottomrule
  \end{tabularx}
  \end{adjustwidth}
\end{table}

\begin{table}[htbp]
  \centering
  \small
  \caption{Complexity, scalability, convergence, and reported
    performance of the surveyed QIEA-MOO methods. Cells marked
    ``not confirmed'' require reading the full paper, not just its
    abstract, before any specific claim is cited (see the methodological
    caveat above).}
  \label{tab:feature-comparison-performance}
  \begin{adjustwidth}{-2cm}{-2cm}
  \centering
  \begin{tabularx}{\linewidth}{@{}p{2.1cm}XXXX@{}}
    \toprule
    Method & Complexity & Scalability & Convergence & Reported
    Performance \\
    \midrule
    QMEA~\citep{kim2006qmea} &
      Not confirmed; direct Pareto ranking without NSGA-II's fast
      non-dominated sort suggests a per-generation cost at or above
      NSGA-II's, not below it &
      Evaluated only on the knapsack instance sizes used in the
      original paper; not confirmed &
      Not confirmed &
      Not confirmed beyond the paper's own framing as an early
      QIEA-MOO proof of concept \\
    QNSGA-II~\citep{guzel2022qnsga2} &
      Inherits NSGA-II's per-generation non-dominated-sorting cost,
      since selection is left unmodified (Section~\ref{sec:taxonomy:selection});
      the superposition-sampling step adds low-order per-Q-bit overhead &
      Not confirmed &
      Reported to converge faster than a plain NSGA-II baseline; the
      magnitude and statistical significance of that claim are not
      independently confirmed here &
      Improved convergence speed relative to NSGA-II (self-reported) \\
    Continuous QIEA-MOO~\citep{olvera2023continuous} &
      Not confirmed &
      Not confirmed &
      Directly compared against NSGA-II on shared continuous benchmark
      problems, per the paper's own framing, but the direction and size
      of any advantage are not confirmed here &
      Not confirmed beyond the paper's framing as a head-to-head
      comparison against NSGA-II \\
    EAH-QNSGA-II~\citep{kumar2026eahqnsga2} &
      Inherits NSGA-II's per-generation cost plus entanglement-crossover
      and local-search overhead; the local-search pass is structurally
      analogous to the hybrid-LAS refinement in this survey's own case
      study (Section~\ref{sec:case-study:mutation}), applied after the
      main GA loop rather than during it &
      Evaluated on a single real-world EV-charging-siting case study;
      scaling behavior beyond that instance not confirmed &
      Not confirmed &
      Reports improved siting trade-offs on its case study relative to
      baseline planning approaches; specific metrics not confirmed here \\
    QIEA-MOO for waste collection~\citep{gilea2026waste} &
      Not confirmed &
      Not confirmed &
      Not confirmed &
      Not confirmed \\
    \bottomrule
  \end{tabularx}
  \end{adjustwidth}
\end{table}

Read across both tables, two patterns stand out even before the
critical analysis proper (Section~\ref{sec:critical-analysis}). First,
every method with a confirmed selection mechanism reuses NSGA-II's
crowding distance unmodified (Table~\ref{tab:feature-comparison-taxonomy}),
so any complexity or scalability claim that does exist is really a
claim about NSGA-II's own well-studied $O(MN^2)$ non-dominated-sorting
cost~\citep{deb2002nsga2} plus a low-order quantum-operator overhead --
none of the surveyed methods has been evaluated against a
selection mechanism whose asymptotic behavior differs enough from
NSGA-II's to matter at scale. Second, reported performance in this
literature is consistently self-comparative (a proposed method against
a classical NSGA-II baseline, run once per paper) rather than
cross-compared against the other QIEA-MOO methods in the same table --
no two of the four confirmed methods above are benchmarked against each
other anywhere in the sources gathered for this survey. Both
observations recur as cross-cutting limitations in
Section~\ref{sec:critical-analysis:cross-cutting}.
